%% toc: P3 28 28.1,28.2,28.3,28.4
\section{Proofs for quantifier elimination, existential closure and model companions}
\label{sec:annex-p03-w03-p03-020-qe}

Two conventions are used throughout. First, when \(h:\mathcal M\to\mathcal N\)
is an embedding, \(\mathcal N\) may be replaced by an isomorphic copy in
which \(h\) is the inclusion: rename each \(h(a)\) as \(a\) and each element
outside the range of \(h\) as a new object not in \(M\), and transport the
interpretations; isomorphic structures satisfy the same formulas at
corresponding tuples by \cref{lem:iso-invariance}. Second, compactness
(\cref{thm:fol-compactness}) is applied to signatures that may be
uncountable, so the Axiom of Choice is assumed wherever it is invoked; the
same assumption covers the construction of chains of length \(\omega\) by
repeated choice.

If \(T\entails\varphi(\bar c)\) for new constants \(\bar c\) not occurring
in \(T\), then \(T\entails\forall\bar x\,\varphi(\bar x)\): given a model of
\(T\) and a tuple, expand the model by interpreting \(\bar c\) as that
tuple. This \emph{generalization on new constants} is used without further
mention.

\begin{annexitem}{annex:p03-w03-p03-020-primitive}{thm:p03-w03-p03-020-primitive}{Syntactic criterion: primitive formulas suffice}
\begin{proof}
A primitive formula is a formula, so elimination gives the condition.
Conversely assume the condition. We take \(\neg\), \(\land\) and \(\exists\)
as primitive, the other connectives and \(\forall\) as abbreviations, and
prove by induction on \(\varphi\) that there is a quantifier-free
\(\varphi^*\) with free variables among those of \(\varphi\) and
\(T\entails\forall\bar x\,(\varphi\leftrightarrow\varphi^*)\).
For atomic \(\varphi\), and for \(\top\) and \(\bot\), put
\(\varphi^*=\varphi\). Put \((\neg\varphi)^*=\neg\varphi^*\) and
\((\varphi_1\land\varphi_2)^*=\varphi_1^*\land\varphi_2^*\); the
equivalences hold because the truth value of a Boolean combination at an
assignment depends only on the truth values of its components there.

Let \(\varphi=\exists y\,\eta\). If two formulas are \(T\)-equivalent, so
are their existential quantifications over \(y\), since in each model of
\(T\) they hold at the same assignments. Hence \(\varphi\) is
\(T\)-equivalent to \(\exists y\,\eta^*\). Write \(\eta^*\) in disjunctive
normal form \(\bigvee_{i<m}\theta_i\), each \(\theta_i\) a finite
conjunction of literals built from the atoms of \(\eta^*\); this is a
tautological equivalence and introduces no variable. If \(m=0\), then
\(\eta^*\) is unsatisfiable and \(\varphi^*=\bot\) works. Otherwise
\(\exists y\bigvee_i\theta_i\) is logically equivalent to
\(\bigvee_i\exists y\,\theta_i\). If \(\theta_i\) is the empty conjunction
\(\top\), then \(\exists y\,\theta_i\) is equivalent to \(\top\) because
structures are nonempty; put \(\psi_i=\top\). Otherwise
\(\exists y\,\theta_i\) is primitive, and the hypothesis, applied with the
tuple of its free variables, gives a quantifier-free \(\psi_i\) with free
variables among those of \(\exists y\,\theta_i\), hence among those of
\(\varphi\), and \(T\)-equivalent to it. Put
\(\varphi^*=\bigvee_i\psi_i\).
\end{proof}
\end{annexitem}

\begin{annexitem}{annex:p03-w03-p03-020-semantic}{thm:p03-w03-p03-020-semantic}{Semantic criterion for one formula}
\begin{proof}
(1) implies (2). Let \(\psi\) be the quantifier-free equivalent. If
\(\mathcal M\entails\varphi(\bar a)\), then \(\mathcal M\entails\psi(\bar a)\),
so \(\psi\in\operatorname{qftp}^{\mathcal M}(\bar a)=
\operatorname{qftp}^{\mathcal N}(\bar b)\), hence
\(\mathcal N\entails\psi(\bar b)\) and \(\mathcal N\entails\varphi(\bar b)\).

(2) implies (1). If \(T\) has no model, \(\top\) is an equivalent. Assume
otherwise and let \(\bar c\) be new constant symbols, one for each variable
of \(\bar x\). Let \(\Gamma\) be the set of quantifier-free \(\psi(\bar x)\)
with \(T\entails\forall\bar x\,(\varphi\to\psi)\). It contains \(\top\)
and is closed under finite conjunctions.

\emph{Claim:} \(T\cup\Gamma(\bar c)\entails\varphi(\bar c)\). Let
\(\mathcal M'\entails T\cup\Gamma(\bar c)\), let \(\mathcal M\) be its
\(\Lang\)-reduct and \(\bar a=\bar c^{\mathcal M'}\). Let \(\Delta\) be the
set of sentences \(\chi(\bar c)\) with \(\chi(\bar x)\) a literal and
\(\mathcal M\entails\chi(\bar a)\). Suppose
\(T\cup\Delta\cup\{\varphi(\bar c)\}\) had no model. By compactness some
\(\chi_1,\ldots,\chi_k\in\Delta\) give
\(T\entails\varphi(\bar c)\to\neg\bigwedge_j\chi_j(\bar c)\), and by
generalization on the new constants
\(T\entails\forall\bar x\,(\varphi\to\neg\bigwedge_j\chi_j)\). So
\(\neg\bigwedge_j\chi_j\in\Gamma\) and
\(\mathcal M\entails\neg\bigwedge_j\chi_j(\bar a)\), which contradicts
\(\chi_j(\bar c)\in\Delta\) for every \(j\) (for \(k=0\) it contradicts
\(\mathcal M\entails\top\)). So there is
\(\mathcal N'\entails T\cup\Delta\cup\{\varphi(\bar c)\}\); let
\(\mathcal N\) be its reduct and \(\bar b=\bar c^{\mathcal N'}\). For each
atomic \(\alpha(\bar x)\), whichever of \(\alpha\) and \(\neg\alpha\) holds
at \(\bar a\) in \(\mathcal M\) lies in \(\Delta\) and therefore holds at
\(\bar b\) in \(\mathcal N\). By induction on Boolean combinations the two
tuples have the same quantifier-free type. Since
\(\mathcal N\entails\varphi(\bar b)\), condition (2), applied with the
roles of the two models exchanged, gives
\(\mathcal M\entails\varphi(\bar a)\). This proves the claim.

By compactness, \(T\cup\{\psi_1(\bar c),\ldots,\psi_k(\bar c)\}\entails
\varphi(\bar c)\) for some \(\psi_j\in\Gamma\). Let \(\psi=\bigwedge_j\psi_j\),
which is \(\top\) if \(k=0\). Then
\(T\entails\forall\bar x\,(\psi\to\varphi)\) by generalization, and
\(T\entails\forall\bar x\,(\varphi\to\psi)\) because \(\psi\in\Gamma\).
\end{proof}
\end{annexitem}

\begin{annexitem}{annex:p03-w03-p03-020-one-point}{thm:p03-w03-p03-020-one-point}{One-point extension test}
\begin{proof}
Assume \(T\) eliminates quantifiers, and let \(\mathcal M,\mathcal N,\bar a,
\bar b,c\) be as in the test. Work in \(\Lang(N)\) with a further new
constant \(e\), and let \(\bar c_{\bar b}\) be the constants naming
\(\bar b\). Let
\[
 \Sigma=\operatorname{ElDiag}(\mathcal N)\cup
 \{\chi(\bar c_{\bar b},e):\chi(\bar x,y)\text{ quantifier-free},\
 \mathcal M\entails\chi(\bar a,c)\}.
\]
A finite subset mentions finitely many such \(\chi_1,\ldots,\chi_k\). Let
\(\theta=\bigwedge_j\chi_j\). Then
\(\mathcal M\entails\exists y\,\theta(\bar a,y)\). By elimination there is a
quantifier-free \(\psi(\bar x)\) that is \(T\)-equivalent to
\(\exists y\,\theta\). So \(\psi\in\operatorname{qftp}^{\mathcal M}(\bar a)=
\operatorname{qftp}^{\mathcal N}(\bar b)\), whence
\(\mathcal N\entails\exists y\,\theta(\bar b,y)\); interpreting \(e\) by a
witness makes \(\mathcal N_N\) a model of the finite subset. By compactness
\(\Sigma\) has a model \(\mathcal N''\). By
\cref{thm:diagram-tests}, item~2, \(n\mapsto c_n^{\mathcal N''}\) is an
elementary embedding of \(\mathcal N\) into the \(\Lang\)-reduct of
\(\mathcal N''\); replacing that reduct by an isomorphic copy
\(\mathcal N'\), we may take the embedding to be the inclusion, so
\(\mathcal N\preccurlyeq\mathcal N'\). Let \(d\) be the element
corresponding to \(e\). Every quantifier-free formula true of
\((\bar a,c)\) in \(\mathcal M\) is true of \((\bar b,d)\) in
\(\mathcal N'\). Since each quantifier-free type contains, for every
quantifier-free \(\chi\), exactly one of \(\chi\) and \(\neg\chi\), the two
types are equal.

Conversely assume the condition. By
\cref{thm:p03-w03-p03-020-primitive} it suffices to treat a primitive
formula \(\exists y\,\theta(\bar x,y)\), and by
\cref{thm:p03-w03-p03-020-semantic} it suffices to verify its condition
(2). Let \(\bar a,\bar b\) have the same quantifier-free type in models
\(\mathcal M,\mathcal N\) of \(T\), and let
\(\mathcal M\entails\theta(\bar a,c)\). Choose \(\mathcal N'\) and \(d\) as
in the condition. Since \(\theta\) is quantifier-free and the types of
\((\bar a,c)\) and \((\bar b,d)\) agree,
\(\mathcal N'\entails\theta(\bar b,d)\), so
\(\mathcal N'\entails\exists y\,\theta(\bar b,y)\), and
\(\mathcal N\preccurlyeq\mathcal N'\) gives
\(\mathcal N\entails\exists y\,\theta(\bar b,y)\).
\end{proof}
\end{annexitem}

\begin{annexitem}{annex:p03-w03-p03-020-qe-consequences}{thm:p03-w03-p03-020-qe-consequences}{What elimination gives}
\begin{proof}
Item 1. Every model of \(T'\) is a model of \(T\), so every
\(T\)-equivalence is a \(T'\)-equivalence.

Item 2. An \(\Lang'\)-formula has the form \(\varphi(\bar x,\bar c)\), where
\(\varphi(\bar x,\bar z)\) is an \(\Lang\)-formula obtained by replacing the
distinct new constants \(\bar c\) by distinct new variables \(\bar z\). Let
\(\psi(\bar x,\bar z)\) be a quantifier-free \(\Lang\)-formula
\(T\)-equivalent to \(\varphi\). Every \(\Lang'\)-model of \(T\) has an
\(\Lang\)-reduct that is a model of \(T\), in which
\(\forall\bar x\,\forall\bar z\,(\varphi\leftrightarrow\psi)\) holds; its
instance at the interpretations of \(\bar c\) shows that
\(\psi(\bar x,\bar c)\) is a quantifier-free \(\Lang'\)-formula
\(T\)-equivalent to \(\varphi(\bar x,\bar c)\).

Item 3. Let \(\mathcal M\subseteq\mathcal N\) be models of \(T\),
\(\varphi(\bar x)\) a formula, \(\bar a\) from \(M\), and \(\psi\) a
quantifier-free \(T\)-equivalent. An inclusion of structures preserves and
reflects atomic formulas at tuples from the smaller structure, hence all
quantifier-free formulas. So
\(\mathcal M\entails\varphi(\bar a)\iff\mathcal M\entails\psi(\bar a)\iff
\mathcal N\entails\psi(\bar a)\iff\mathcal N\entails\varphi(\bar a)\). Thus
\(\mathcal M\preccurlyeq\mathcal N\), and
\cref{lem:model-completeness-reduction} gives model completeness.

Item 4. Let \(\sigma\) be an \(\Lang(A)\)-sentence and write it as
\(\varphi(\bar c_{\bar a})\), where \(\varphi(\bar x)\) is an
\(\Lang\)-formula and \(\bar c_{\bar a}\) are the distinct parameter
constants occurring in \(\sigma\). Let \(\psi\) be a quantifier-free
\(T\)-equivalent of \(\varphi\). Every model of
\(T\cup\operatorname{Diag}(\mathcal A)\) has an \(\Lang\)-reduct that is a
model of \(T\), so the theory entails
\(\sigma\leftrightarrow\psi(\bar c_{\bar a})\). The sentence
\(\psi(\bar c_{\bar a})\) is a Boolean combination of atomic
\(\Lang(A)\)-sentences, and for each atomic \(\Lang(A)\)-sentence
\(\beta\), either \(\beta\) or \(\neg\beta\) belongs to
\(\operatorname{Diag}(\mathcal A)\). By induction on the Boolean
combination, \(\operatorname{Diag}(\mathcal A)\) entails
\(\psi(\bar c_{\bar a})\) or its negation, and so the theory entails
\(\sigma\) or \(\neg\sigma\). By \cref{thm:diagram-tests}, item~1, the
theory has a model exactly when \(\mathcal A\) embeds in a model of \(T\).

Item 5. Let \(\sigma\) be a sentence and \(\psi\) a quantifier-free
\(T\)-equivalent sentence, a Boolean combination of atomic sentences (and
\(\top,\bot\)). Any two models of \(T\) agree on the atomic sentences,
hence on \(\psi\), hence on \(\sigma\). Fix a model \(\mathcal M\). If
\(\mathcal M\entails\sigma\), every model satisfies \(\sigma\) and
\(T\entails\sigma\); otherwise \(T\entails\neg\sigma\). Semantic and
deductive completeness agree by \cref{def:theory-completeness}.
\end{proof}
\end{annexitem}

\begin{annexitem}{annex:p03-w03-p03-020-existential-preservation}{thm:p03-w03-p03-020-existential-preservation}{Existential normal form from preservation}
\begin{proof}
If \(\varphi\) is \(T\)-equivalent to \(\exists\bar y\,\theta\) with
\(\theta\) quantifier-free and \(\mathcal M\entails\theta(\bar a,\bar e)\),
then \(\mathcal N\entails\theta(\bar a,\bar e)\) because the inclusion
preserves quantifier-free formulas, so truth passes upward.

Conversely suppose truth of \(\varphi\) passes upward along inclusions of
models of \(T\); equivalently, truth of \(\neg\varphi\) passes downward.
Let \(\Psi\) be the set of universal \(\psi(\bar x)\) with
\(T\entails\forall\bar x\,(\neg\varphi\to\psi)\), and let \(\bar c\) be
new constants.

\emph{Claim:} \(T\cup\Psi(\bar c)\entails\neg\varphi(\bar c)\). Let
\(\mathcal N'\entails T\cup\Psi(\bar c)\), with reduct \(\mathcal N\) and
\(\bar a=\bar c^{\mathcal N'}\). Consider the \(\Lang(N)\cup\{\bar c\}\)-theory
\[
 \Sigma=T\cup\operatorname{Diag}(\mathcal N)\cup
 \{c_i=c_{a_i}:i\le n\}\cup\{\neg\varphi(\bar c)\}.
\]
Suppose \(\Sigma\) had no model. By compactness there is a finite
conjunction \(\delta\) of members of \(\operatorname{Diag}(\mathcal N)\),
whose parameter constants we may assume include all \(c_{a_i}\), such that
\(T\cup\{\delta\}\cup\{c_i=c_{a_i}\}\entails\varphi(\bar c)\). Replace each
parameter constant \(c_e\) in \(\delta\) by a new variable \(z_e\), distinct
constants by distinct variables, obtaining \(\delta(\bar z)\). Generalization
gives
\(T\entails\forall\bar x\,\forall\bar z\,
((\delta(\bar z)\land\bigwedge_i x_i=z_{a_i})\to\varphi(\bar x))\), hence
\(T\entails\forall\bar x\,(\neg\varphi\to
\forall\bar z\,\neg(\delta(\bar z)\land\bigwedge_i x_i=z_{a_i}))\). The
consequent is universal, so it lies in \(\Psi\), and
\(\mathcal N\) satisfies it at \(\bar a\). But the elements named in
\(\delta\) satisfy \(\delta(\bar z)\land\bigwedge_i a_i=z_{a_i}\) in
\(\mathcal N\), a contradiction. So \(\Sigma\) has a model \(\mathcal P'\).
By \cref{thm:diagram-tests}, item~1, \(e\mapsto c_e^{\mathcal P'}\) embeds
\(\mathcal N\) in the \(\Lang\)-reduct of \(\mathcal P'\); after the
renaming convention this is an inclusion \(\mathcal N\subseteq\mathcal P\)
of models of \(T\) with \(\bar c^{\mathcal P'}=\bar a\) and
\(\mathcal P\entails\neg\varphi(\bar a)\). By downward passage,
\(\mathcal N\entails\neg\varphi(\bar a)\), proving the claim.

By compactness \(T\cup\{\psi_1(\bar c),\ldots,\psi_k(\bar c)\}\entails
\neg\varphi(\bar c)\) for some \(\psi_j\in\Psi\). After renaming bound
variables apart, \(\bigwedge_j\psi_j\) is logically equivalent to a single
universal formula \(\psi\) (with \(\psi=\top\) when \(k=0\)). Then
\(T\entails\forall\bar x\,(\psi\to\neg\varphi)\), and
\(T\entails\forall\bar x\,(\neg\varphi\to\psi)\) because each
\(\psi_j\in\Psi\). So \(\neg\varphi\) is \(T\)-equivalent to
\(\psi=\forall\bar y\,\eta\), and \(\varphi\) to the existential formula
\(\exists\bar y\,\neg\eta\).

The dual statement follows by applying this one to \(\neg\varphi\): the
formula \(\varphi\) is \(T\)-equivalent to a universal formula iff
\(\neg\varphi\) is \(T\)-equivalent to an existential one, iff truth of
\(\neg\varphi\) passes upward, iff truth of \(\varphi\) passes downward.
\end{proof}
\end{annexitem}

\begin{annexitem}{annex:p03-w03-p03-020-robinson}{thm:p03-w03-p03-020-robinson}{Robinson's test, with its diagram and normal-form bridges}
\begin{proof}
(1) implies (2). For models \(\mathcal M\subseteq\mathcal N\) of \(T\) the
inclusion is an embedding between models, hence elementary; in particular
existential \(\Lang(M)\)-sentences true in \(\mathcal N_M\) are true in
\(\mathcal M_M\).

(2) implies (3). By induction on formulas built from atoms with \(\neg\),
\(\land\) and \(\exists\), every formula is \(T\)-equivalent to an
existential formula. Atoms are existential with an empty block. If
\(\varphi_i\) is \(T\)-equivalent to \(\exists\bar y_i\,\theta_i\) for
\(i=1,2\), with the blocks renamed apart and not free in the other formula,
then \(\varphi_1\land\varphi_2\) is \(T\)-equivalent to
\(\exists\bar y_1\bar y_2\,(\theta_1\land\theta_2)\), and
\(\exists z\,\varphi_1\) to \(\exists z\,\exists\bar y_1\,\theta_1\). For
negation, \(\neg\varphi_1\) is \(T\)-equivalent to the universal formula
\(\chi=\forall\bar y_1\,\neg\theta_1\). Truth of \(\chi\) passes upward
along inclusions of models of \(T\): if \(\mathcal M\subseteq\mathcal N\),
\(\mathcal M\entails\chi(\bar a)\) and
\(\mathcal N\entails\exists\bar y_1\,\theta_1(\bar a,\bar y_1)\), then by
(2) the existential \(\Lang(M)\)-sentence
\(\exists\bar y_1\,\theta_1(\bar a,\bar y_1)\) holds in \(\mathcal M\),
contrary to \(\mathcal M\entails\chi(\bar a)\). By
\cref{thm:p03-w03-p03-020-existential-preservation}, \(\chi\) and hence
\(\neg\varphi_1\) is \(T\)-equivalent to an existential formula.

(3) and (4) are equivalent. If every formula has an existential
\(T\)-equivalent, apply this to \(\neg\varphi\): from
\(\neg\varphi\leftrightarrow\exists\bar y\,\theta\) we get
\(\varphi\leftrightarrow\forall\bar y\,\neg\theta\). The converse is the
same argument with the quantifiers exchanged.

(3) implies (1). Let \(\mathcal M\subseteq\mathcal N\) be models of \(T\)
and \(\bar a\) from \(M\). Let \(\psi\) and \(\psi'\) be existential
\(T\)-equivalents of \(\varphi\) and \(\neg\varphi\). If
\(\mathcal M\entails\varphi(\bar a)\), then
\(\mathcal M\entails\psi(\bar a)\), so \(\mathcal N\entails\psi(\bar a)\)
by the easy half of \cref{thm:p03-w03-p03-020-existential-preservation},
and \(\mathcal N\entails\varphi(\bar a)\). The same argument with
\(\psi'\) handles \(\neg\varphi\). So \(\mathcal M\preccurlyeq\mathcal N\),
and \cref{lem:model-completeness-reduction} gives (1).

(1) implies (5). Let \(\mathcal M\entails T\) and let
\(\mathcal P'\entails T\cup\operatorname{Diag}(\mathcal M)\). By
\cref{thm:diagram-tests}, item~1, \(h(a)=c_a^{\mathcal P'}\) is an
embedding of \(\mathcal M\) into the reduct \(\mathcal P\), a model of
\(T\); by (1) it is elementary. Hence for every \(\Lang(M)\)-sentence
\(\varphi(\bar c_{\bar a})\) true in \(\mathcal M_M\),
\(\mathcal P\entails\varphi(h\bar a)\), that is,
\(\mathcal P'\entails\varphi(\bar c_{\bar a})\). So
\(T\cup\operatorname{Diag}(\mathcal M)\) entails
\(\operatorname{ElDiag}(\mathcal M)\), which contains every
\(\Lang(M)\)-sentence or its negation; thus the theory is complete. The two
forms of (5) are equivalent: a complete theory with the model
\(\mathcal M_M\) entails every sentence true in \(\mathcal M_M\), and a
theory entailing the complete set \(\operatorname{ElDiag}(\mathcal M)\) is
complete.

(5) implies (1). Let \(\mathcal M\subseteq\mathcal N\) be models of \(T\).
The expansion \(\mathcal N_M\) with \(c_a\) interpreted as \(a\) satisfies
\(\operatorname{Diag}(\mathcal M)\), because the inclusion is an embedding,
so it satisfies \(\operatorname{ElDiag}(\mathcal M)\). Thus every formula
true of a tuple from \(M\) in \(\mathcal M\) is true of it in
\(\mathcal N\), so \(\mathcal M\preccurlyeq\mathcal N\), and
\cref{lem:model-completeness-reduction} gives (1).
\end{proof}
\end{annexitem}

\begin{annexitem}{annex:p03-w03-p03-020-ec-elementary}{thm:p03-w03-p03-020-ec-elementary}{Existential closure and elementary extensions}
\begin{proof}
Since \(A\subseteq B\), the signature \(\Lang(A)\) is part of \(\Lang(B)\).
Let \(\Sigma=\operatorname{ElDiag}(\mathcal A)\cup\operatorname{Diag}(\mathcal B)\).
Suppose \(\Sigma\) had no model. By compactness there is a finite
conjunction \(\delta(\bar c_{\bar a},\bar c_{\bar e})\) of members of
\(\operatorname{Diag}(\mathcal B)\), with \(\bar a\) from \(A\) and
\(\bar e\) distinct elements of \(B\setminus A\), such that
\(\operatorname{ElDiag}(\mathcal A)\entails\neg\delta(\bar c_{\bar a},
\bar c_{\bar e})\). The constants \(\bar c_{\bar e}\) do not occur in
\(\operatorname{ElDiag}(\mathcal A)\), so by generalization it entails
\(\forall\bar z\,\neg\delta(\bar c_{\bar a},\bar z)\), which is therefore
true in \(\mathcal A_A\). But \(\mathcal B\entails\delta(\bar a,\bar e)\), so
the existential \(\Lang(A)\)-sentence
\(\exists\bar z\,\delta(\bar c_{\bar a},\bar z)\) is true in
\(\mathcal B_A\) and, by existential closure, in \(\mathcal A_A\): a
contradiction. Let \(\mathcal A''\) be a model of \(\Sigma\). By
\cref{thm:diagram-tests}, \(a\mapsto c_a^{\mathcal A''}\) is an elementary
embedding of \(\mathcal A\) and \(b\mapsto c_b^{\mathcal A''}\) an embedding of
\(\mathcal B\) into the \(\Lang\)-reduct, and the two maps agree on \(A\).
After renaming, the first is an inclusion \(\mathcal A\preccurlyeq
\mathcal A'\), and the second is then the identity on \(A\).
\end{proof}
\end{annexitem}

\begin{annexitem}{annex:p03-w03-p03-020-ae}{thm:p03-w03-p03-020-ae}{Model complete theories are \(\forall\exists\)-axiomatized}
\begin{proof}
Every model of \(T\) satisfies \(T_{\forall\exists}\). Conversely let
\(\mathcal A\entails T_{\forall\exists}\).

\emph{Step.} There is a model \(\mathcal B\) of \(T\) with
\(\mathcal A\subseteq\mathcal B\) and \(\mathcal A\) existentially closed in
\(\mathcal B\). Let \(U\) be the set of universal \(\Lang(A)\)-sentences
true in \(\mathcal A_A\) (the negations of the false existential ones), and
\(\Sigma=T\cup\operatorname{Diag}(\mathcal A)\cup U\). Suppose \(\Sigma\)
had no model. By compactness some finite conjunction \(\delta\) of members
of the diagram and some \(u\in U\) (a finite conjunction of members of
\(U\) is equivalent to a universal sentence after renaming bound variables)
satisfy \(T\entails\neg(\delta\land u)\). Write
\(u=\forall\bar w\,\eta(\bar c_{\bar a},\bar w)\) with \(\eta\)
quantifier-free and replace the distinct parameter constants
\(\bar c_{\bar a}\) of \(\delta\) and \(u\) by distinct variables \(\bar x\).
Generalization gives
\(T\entails\forall\bar x\,\exists\bar w\,(\neg\delta(\bar x)\lor
\neg\eta(\bar x,\bar w))\), a sentence of \(T_{\forall\exists}\). It
therefore holds in \(\mathcal A\), and at \(\bar a\) it contradicts
\(\mathcal A\entails\delta(\bar a)\land\forall\bar w\,\eta(\bar a,\bar w)\).
So \(\Sigma\) has a model; by \cref{thm:diagram-tests} and renaming its
reduct is a model \(\mathcal B\supseteq\mathcal A\) of \(T\). An
existential \(\Lang(A)\)-sentence false in \(\mathcal A_A\) has its
negation in \(U\), so it is false in \(\mathcal B_A\).

\emph{Chain.} By the step and
\cref{thm:p03-w03-p03-020-ec-elementary} there are
\(\mathcal B_0\entails T\) and \(\mathcal A_1\) with
\(\mathcal A_0=\mathcal A\subseteq\mathcal B_0\subseteq\mathcal A_1\) and
\(\mathcal A_0\preccurlyeq\mathcal A_1\). Then \(\mathcal A_1\) satisfies
the same sentences as \(\mathcal A_0\), in particular
\(T_{\forall\exists}\), so the construction repeats. Choosing at each stage
gives \(\mathcal A_0\subseteq\mathcal B_0\subseteq\mathcal A_1\subseteq
\mathcal B_1\subseteq\cdots\) with \(\mathcal A_i\preccurlyeq\mathcal A_{i+1}\)
and every \(\mathcal B_i\) a model of \(T\). By model completeness the
inclusions \(\mathcal B_i\subseteq\mathcal B_{i+1}\) are elementary. The
sequences \((\mathcal A_i)\) and \((\mathcal B_i)\) interleave, so they have
the same union \(\mathcal U\). By \cref{thm:elementary-chain},
\(\mathcal A_0\preccurlyeq\mathcal U\) and
\(\mathcal B_0\preccurlyeq\mathcal U\). Hence \(\mathcal A\) satisfies the
same sentences as \(\mathcal B_0\), which is a model of \(T\).
\end{proof}
\end{annexitem}

\begin{annexitem}{annex:p03-w03-p03-020-companion-embedding}{thm:p03-w03-p03-020-companion-embedding}{Companions by mutual embedding}
\begin{proof}
Assume \(T_\forall=T^*_\forall\) and let \(\mathcal M\entails T\). Suppose
\(T^*\cup\operatorname{Diag}(\mathcal M)\) had no model. By compactness
some finite conjunction \(\delta(\bar c_{\bar a})\) of diagram sentences, with
\(\bar a\) distinct, satisfies \(T^*\entails\neg\delta(\bar c_{\bar a})\), and
by generalization \(T^*\entails\forall\bar x\,\neg\delta(\bar x)\). This
universal sentence lies in \(T^*_\forall=T_\forall\), so it holds in
\(\mathcal M\), contradicting \(\mathcal M\entails\delta(\bar a)\). A model
of \(T^*\cup\operatorname{Diag}(\mathcal M)\) gives an embedding of
\(\mathcal M\) into a model of \(T^*\) by \cref{thm:diagram-tests}. The
other direction is symmetric.

Conversely assume the mutual embeddings and let
\(\sigma=\forall\bar x\,\theta\in T^*_\forall\) with \(\theta\)
quantifier-free. Let \(\mathcal M\entails T\) and
\(h:\mathcal M\to\mathcal N\entails T^*\). For \(\bar a\) from \(M\),
\(\mathcal N\entails\theta(h\bar a)\), and an embedding reflects
quantifier-free formulas, so \(\mathcal M\entails\theta(\bar a)\). Thus
\(T\entails\sigma\) and \(\sigma\in T_\forall\). By symmetry the two sets
are equal.
\end{proof}
\end{annexitem}

\begin{annexitem}{annex:p03-w03-p03-020-companion-unique}{thm:p03-w03-p03-020-companion-unique}{Uniqueness of model companions}
\begin{proof}
Both \(T_1^*\) and \(T_2^*\) have the universal consequences of \(T\), so
\(T_{1\forall}^*=T_{2\forall}^*\) and
\cref{thm:p03-w03-p03-020-companion-embedding} applies to the pair. Let
\(\mathcal M_0\entails T_1^*\). Using the lemma alternately and renaming
each time, build
\(\mathcal M_0\subseteq\mathcal N_0\subseteq\mathcal M_1\subseteq
\mathcal N_1\subseteq\cdots\) with \(\mathcal M_i\entails T_1^*\) and
\(\mathcal N_i\entails T_2^*\). Both theories are model complete, so
\(\mathcal M_i\preccurlyeq\mathcal M_{i+1}\) and
\(\mathcal N_i\preccurlyeq\mathcal N_{i+1}\). The two chains have the same
union \(\mathcal U\), and \cref{thm:elementary-chain} gives
\(\mathcal M_0\preccurlyeq\mathcal U\) and
\(\mathcal N_0\preccurlyeq\mathcal U\). So \(\mathcal M_0\) satisfies the
same sentences as \(\mathcal N_0\), a model of \(T_2^*\). By symmetry the
two theories have the same models. A model completion is by definition a
model companion, so the same holds for model completions.
\end{proof}
\end{annexitem}

\begin{annexitem}{annex:p03-w03-p03-020-companion-ec}{thm:p03-w03-p03-020-companion-ec}{Models of a model companion are the existentially closed models}
\begin{proof}
Item 1. Let \(\mathcal M\entails T^*\) and
\(\mathcal M\subseteq\mathcal N\entails T_\forall=T^*_\forall\). The first
half of the proof of \cref{thm:p03-w03-p03-020-companion-embedding} used
only that \(\mathcal N\) satisfies the universal consequences of \(T^*\),
so it embeds \(\mathcal N\) in a model of \(T^*\); after renaming,
\(\mathcal N\subseteq\mathcal N'\entails T^*\). Model completeness gives
\(\mathcal M\preccurlyeq\mathcal N'\). An existential
\(\Lang(M)\)-sentence true in \(\mathcal N\) is true in \(\mathcal N'\),
because its witnesses and the quantifier-free matrix pass up the inclusion,
and then it is true in \(\mathcal M\). Models of \(T\) are models of
\(T_\forall\).

Item 2. Assume \(T\) has the same models as \(T_{\forall\exists}\).
Let \(\mathcal M\entails T^*\). By
\cref{thm:p03-w03-p03-020-companion-embedding} and renaming,
\(\mathcal M\subseteq\mathcal N\) for some \(\mathcal N\entails T\), and
\(\mathcal M\) is existentially closed in \(\mathcal N\) by item~1. For
\(\forall\bar x\,\exists\bar y\,\theta\in T_{\forall\exists}\) and \(\bar a\)
from \(M\), \(\mathcal N\entails\exists\bar y\,\theta(\bar a,\bar y)\), so
\(\mathcal M\entails\exists\bar y\,\theta(\bar a,\bar y)\). Thus
\(\mathcal M\entails T\), and item~1 shows it is existentially closed among
models of \(T\).

Conversely let \(\mathcal M\entails T\) be existentially closed among models
of \(T\). By the lemma and renaming,
\(\mathcal M\subseteq\mathcal N\subseteq\mathcal P\) with
\(\mathcal N\entails T^*\) and \(\mathcal P\entails T\). An existential
\(\Lang(M)\)-sentence true in \(\mathcal N\) is true in \(\mathcal P\) and
hence in \(\mathcal M\). By \cref{thm:p03-w03-p03-020-ae}, \(T^*\) has
the same models as \(T^*_{\forall\exists}\); for each of its sentences
\(\forall\bar x\,\exists\bar y\,\theta\) and \(\bar a\) from \(M\), the
sentence \(\exists\bar y\,\theta(\bar a,\bar y)\) holds in \(\mathcal N\) and
so in \(\mathcal M\). Hence \(\mathcal M\entails T^*\).
\end{proof}
\end{annexitem}

\begin{annexitem}{annex:p03-w03-p03-020-qe-completion}{thm:p03-w03-p03-020-qe-completion}{A companion with elimination is a model completion}
\begin{proof}
By \cref{thm:p03-w03-p03-020-qe-consequences}, item~3, \(T^*\) is model
complete, so it is a model companion of \(T\) by
\cref{def:p03-model-companion}. Let \(\mathcal M\entails T\). By
\cref{thm:p03-w03-p03-020-companion-embedding}, \(\mathcal M\) embeds in a
model of \(T^*\), so \(T^*\cup\operatorname{Diag}(\mathcal M)\) is
consistent by \cref{thm:diagram-tests}, and it is complete by
\cref{thm:p03-w03-p03-020-qe-consequences}, item~4. So \(T^*\) is a model
completion of \(T\), and the only one up to models by
\cref{thm:p03-w03-p03-020-companion-unique}.
\end{proof}
\end{annexitem}

\begin{annexitem}{annex:p03-w03-p03-020-bf-preserve}{thm:p03-w03-p03-020-bf-preserve}{Back-and-forth systems preserve all formulas}
\begin{proof}
Take \(\neg\), \(\land\) and \(\exists\) as primitive. We prove the
statement for all \(p\in I\) and all tuples from the domain of \(p\)
simultaneously, by induction on \(\varphi\). For atomic \(\varphi\) it is
the definition of a finite partial isomorphism, and the Boolean steps
follow from the induction hypothesis. Let
\(\varphi=\exists y\,\eta(\bar x,y)\). If
\(\mathcal M\entails\eta(\bar a,c)\), the forth property gives
\(q\in I\) extending \(p\) with \(c\) in its domain. The tuple
\((\bar a,c)\) lies in the domain of \(q\), and \(q\bar a=p\bar a\), so the
induction hypothesis for \(\eta\) at \(q\) gives
\(\mathcal N\entails\eta(p\bar a,q(c))\), whence
\(\mathcal N\entails\varphi(p\bar a)\). If
\(\mathcal N\entails\eta(p\bar a,d)\), the back property gives \(q\in I\)
extending \(p\) with \(d=q(c)\) for some \(c\) in its domain, and the
induction hypothesis gives \(\mathcal M\entails\eta(\bar a,c)\).
\end{proof}
\end{annexitem}

\begin{annexitem}{annex:p03-w03-p03-020-bf-criterion}{thm:p03-w03-p03-020-bf-criterion}{Back-and-forth criterion for elimination}
\begin{proof}
Let \(\mathcal M,\mathcal N\entails T\) and let \(I\) be the set of all
finite partial isomorphisms from \(\mathcal M\) to \(\mathcal N\). By
hypothesis every member of \(I\) has the forth and back extensions inside
\(I\). We verify condition (2) of
\cref{thm:p03-w03-p03-020-semantic} for every formula. Let
\(\bar a\in M^n\) and \(\bar b\in N^n\) have the same quantifier-free type.
Put \(p(a_i)=b_i\). This is well defined and injective, since
\(a_i=a_j\) iff \(x_i=x_j\) holds at \(\bar a\) iff \(b_i=b_j\). Every
atomic formula about a tuple from \(\{a_1,\ldots,a_n\}\) is an atomic
formula \(\alpha(\bar x)\) about \(\bar a\) after renaming variables, so
\(p\in I\). Thus \(I\) is a back-and-forth system, and
\cref{thm:p03-w03-p03-020-bf-preserve} gives
\(\mathcal M\entails\varphi(\bar a)\iff\mathcal N\entails\varphi(\bar b)\).
Hence every formula has a quantifier-free \(T\)-equivalent. If the models
satisfy the same atomic sentences, the empty map is in \(I\), and the
theorem applied to the empty tuple shows that they satisfy the same
sentences.
\end{proof}
\end{annexitem}

\begin{annexitem}{annex:p03-w03-p03-020-dlo-bf}{thm:p03-w03-p03-020-dlo-bf}{Elimination for dense linear orders by back-and-forth}
\begin{proof}
In \(\{<\}\) the atomic formulas are \(u<v\) and \(u=v\) for variables
\(u,v\). So a finite map \(p\) is a finite partial isomorphism iff it is
injective and \(a<a'\iff p(a)<p(a')\) on its domain \(D\).

\emph{Forth.} Let \(c\in M\). If \(c\in D\), take \(q=p\). Otherwise let
\(L=\{a\in D:a<c\}\) and \(U=\{a\in D:c<a\}\); as the order is linear and
\(c\notin D\), \(D=L\cup U\). Choose \(d\in N\) as follows.
If \(L\) and \(U\) are nonempty, let \(l=\max L\) and \(u=\min U\), which
exist because they are finite nonempty subsets of a linear order; then
\(l<c<u\), so \(p(l)<p(u)\), and density gives \(d\) with
\(p(l)<d<p(u)\). If \(L\ne\emptyset=U\), the absence of a greatest element
gives \(d>p(\max L)\). If \(L=\emptyset\ne U\), the absence of a least
element gives \(d<p(\min U)\). If \(D=\emptyset\), any \(d\in N\) will do,
as \(N\) is nonempty. In each case, for \(a\in L\) we have
\(p(a)\le p(l)<d\), and for \(a\in U\) we have \(d<p(u)\le p(a)\). Hence
\(d\notin p[D]\), and \(q=p\cup\{(c,d)\}\) is injective and matches every
comparison between \(c\) and a member of \(D\) with the corresponding
comparison between \(d\) and its image. So \(q\) is a finite partial
isomorphism.

\emph{Back.} The inverse of a finite partial isomorphism from
\(\mathcal M\) to \(\mathcal N\) is one from \(\mathcal N\) to
\(\mathcal M\), and \(\mathcal N\) is also a model of \(\mathsf{DLO}\), so
the forth step for \(p^{-1}\) gives the back step for \(p\).

By \cref{thm:p03-w03-p03-020-bf-criterion}, \(\mathsf{DLO}\) eliminates
quantifiers. The signature has no constant or nullary relation symbols, so
there are no atomic sentences and any two models are elementarily
equivalent by the same corollary. The theory is consistent, since
\((\mathbb Q;<)\) is a nonempty linear order with
\(r<\tfrac{r+s}2<s\) for \(r<s\) and \(r-1<r<r+1\). Hence it is complete,
and it is model complete by
\cref{thm:p03-w03-p03-020-qe-consequences}, item~3.
\end{proof}
\end{annexitem}

\begin{annexitem}{annex:p03-w03-p03-020-dense-variants}{thm:p03-w03-p03-020-dense-variants}{Elimination for each endpoint pattern}
\begin{proof}
Let \(\mathcal M,\mathcal N\entails\mathsf{DLO}^{E}\). Write
\(e^{\mathcal M}\) for the interpretation of a constant \(e\in E\). The
atomic \(\Lang_E\)-formulas are \(u<v\) and \(u=v\) with \(u,v\) variables
or constants.

\emph{Adding the constants.} Let \(p\) be a finite partial isomorphism
with domain \(D\), and put
\(p'=p\cup\{(e^{\mathcal M},e^{\mathcal N}):e\in E\}\). If
\(e^{\mathcal M}=a\in D\), the atom \(x=e\) holds at \(a\), so
\(p(a)=e^{\mathcal N}\); thus \(p'\) is a function. Comparisons between a
member of \(D\) and a constant are atomic formulas about that member, so
they are preserved and reflected by \(p\). If \(E=\{\mathsf b,\mathsf t\}\),
both models satisfy \(\mathsf b<\mathsf t\), because each has at least two
elements, a least element \(\mathsf b\) and a greatest element
\(\mathsf t\). So \(p'\) is injective and preserves and reflects the order
on its domain \(D'=D\cup\{e^{\mathcal M}:e\in E\}\), and every atomic
formula with parameters in \(D'\) compares members of \(D'\). Hence \(p'\)
is a finite partial isomorphism.

\emph{Forth.} Let \(c\in M\setminus D'\), and define \(L,U\subseteq D'\) as
in the proof of \cref{thm:p03-w03-p03-020-dlo-bf}. If \(U=\emptyset\), then
\(\mathsf t\notin E\): otherwise \(\mathsf t^{\mathcal M}\in D'\), and
\(c<\mathsf t^{\mathcal M}\) because \(c\ne\mathsf t^{\mathcal M}\) and
\(\mathsf t^{\mathcal M}\) is greatest, so \(\mathsf t^{\mathcal M}\in U\).
So \(\mathcal N\) has no greatest element. Symmetrically, if
\(L=\emptyset\), then \(\mathsf b\notin E\) and \(\mathcal N\) has no least
element. If \(L=U=\emptyset\), then \(D'=\emptyset\). The choice of \(d\)
and the verification that \(p'\cup\{(c,d)\}\) is a finite partial
isomorphism are now word for word those of
\cref{thm:p03-w03-p03-020-dlo-bf}, using density for nonempty \(L\) and
\(U\), the absence of the relevant endpoint in \(\mathcal N\) when one of
them is empty, and nonemptiness when both are. Atomic formulas involving
constants compare \(c\) or \(d\) with members of \(D'\) or their images,
which are covered.

\emph{Back} follows by applying the forth step to the inverse map, as
before. By \cref{thm:p03-w03-p03-020-bf-criterion},
\(\mathsf{DLO}^{E}\) eliminates quantifiers. Its atomic sentences compare
the constants of \(E\), and the axioms decide each of them: \(e=e\) and
\(\neg(e<e)\) hold, and if \(E=\{\mathsf b,\mathsf t\}\) then
\(\mathsf b<\mathsf t\), \(\neg(\mathsf t<\mathsf b)\) and
\(\mathsf b\ne\mathsf t\). The theory is consistent: \(\mathbb Q\),
\(\mathbb Q\cap[0,\infty)\), \(\mathbb Q\cap(-\infty,0]\) and
\(\mathbb Q\cap[0,1]\), with the evident constants, are models for
\(E=\emptyset,\{\mathsf b\},\{\mathsf t\},\{\mathsf b,\mathsf t\}\); in each,
midpoints give density and \(r\pm1\) shows the absence of a missing
endpoint. So \(\mathsf{DLO}^{E}\) is complete by
\cref{thm:p03-w03-p03-020-qe-consequences}, item~5, and model complete by
item~3.

\emph{The reducts in the warning.} Let \(T\) be the set of
\(\{<\}\)-sentences entailed by \(\mathsf{DLO}^{E}\). For each
\(\{<\}\)-sentence \(\sigma\), the complete theory \(\mathsf{DLO}^{E}\)
entails \(\sigma\) or \(\neg\sigma\), and so does \(T\); hence \(T\) is
complete. A model of \(T\) satisfies the \(\{<\}\)-sentences expressing
linearity, density, at least two elements and the prescribed existence or
absence of each endpoint, so naming its endpoints yields a model of
\(\mathsf{DLO}^{E}\); thus the models of \(T\) are the reducts of the
models of \(\mathsf{DLO}^{E}\).

\emph{Model completion of linear orders.} Every model of \(\mathsf{DLO}\) is a
strict linear order. Conversely let \((L;<)\) be a nonempty strict linear
order and order \(L\times\mathbb Q\) lexicographically:
\((a,r)<(b,s)\) iff \(a<b\), or \(a=b\) and \(r<s\). This is a strict
linear order. If \((a,r)<(b,s)\), then \((a,\tfrac{r+s}2)\) lies between
them when \(a=b\), and \((a,r+1)\) when \(a<b\); and
\((a,r-1)<(a,r)<(a,r+1)\). So \(L\times\mathbb Q\) is a model of
\(\mathsf{DLO}\), and \(a\mapsto(a,0)\) preserves and reflects \(<\) and
\(=\), so it is an embedding. By
\cref{thm:p03-w03-p03-020-companion-embedding} the two theories have the
same universal consequences, and \cref{thm:p03-w03-p03-020-qe-completion}
shows that \(\mathsf{DLO}\) is the model completion of the theory of
linear orders.
\end{proof}
\end{annexitem}

\begin{annexitem}{annex:p03-w03-p03-020-marked-point}{thm:p03-w03-p03-020-marked-point}{A model complete theory without elimination}
\begin{proof}
Let \(c\) be a new constant. By
\cref{thm:p03-w03-p03-020-dlo-bf} and
\cref{thm:p03-w03-p03-020-qe-consequences}, item~2, \(\mathsf{DLO}\) read
in \(\{<,c\}\) eliminates quantifiers, so it is model complete there.

Let \(\mathcal M\subseteq\mathcal N\) be models of \(T_Z\). Since the
inclusion preserves and reflects \(Z\), \(Z^{\mathcal M}=Z^{\mathcal N}\cap M\).
As \(\mathcal M\) has an element \(z\) in \(Z^{\mathcal M}\), it lies in
\(Z^{\mathcal N}\), and \(Z^{\mathcal N}=\{z\}\) because \(\mathcal N\) has
exactly one such element. Expand both by interpreting \(c\) as \(z\); then
\(\mathcal M_c\subseteq\mathcal N_c\) are models of \(\mathsf{DLO}\) in
\(\{<,c\}\), so \(\mathcal M_c\preccurlyeq\mathcal N_c\). For a
\(\{<,Z\}\)-formula \(\varphi\), let \(\varphi^c\) replace every atom
\(Z(t)\) by \(t=c\). In each of the two models, \(Z\) holds exactly at the
interpretation of \(c\), so by induction on \(\varphi\) a tuple satisfies
\(\varphi\) in the model iff it satisfies \(\varphi^c\) in the expansion.
Hence \(\mathcal M\preccurlyeq\mathcal N\), and
\cref{lem:model-completeness-reduction} gives model completeness.

The structure \((\mathbb Q;<,\{0\})\) is a model of \(T_Z\). The
quantifier-free \(\{<,Z\}\)-formulas in one variable \(x\) are Boolean
combinations of \(x<x\), \(x=x\) and \(Z(x)\), and \(-1\) and \(1\) satisfy
the same ones. But \(\exists y\,(Z(y)\land x<y)\) holds at \(-1\) and fails
at \(1\). By \cref{thm:p03-w03-p03-020-semantic}, this formula has no
quantifier-free \(T_Z\)-equivalent.
\end{proof}
\end{annexitem}

\begin{annexitem}{annex:p03-w03-p03-020-field-algebra}{thm:p03-w03-p03-020-field-algebra}{Field extensions used in the elimination}
\begin{proof}
Item 1. \(E[X]\) is a principal ideal domain with a division algorithm. If
\(h\notin(g)\), then a greatest common divisor of \(g\) and \(h\) divides the
irreducible \(g\) and is not an associate of \(g\), since \(g\nmid h\); so it
is a unit, and \(uh+vg=1\) for some \(u,v\). Thus \(h\) is invertible modulo
\(g\), \((g)\) is maximal, and \(E[X]/(g)\) is a field. The composite
\(E\to E[X]\to E[X]/(g)\) is a ring homomorphism of fields sending \(1\) to
\(1\ne0\), so it is injective; identify \(E\) with its image. If \(\xi\) is
the class of \(X\), then \(g(\xi)\) is the class of \(g\), which is \(0\).
For finitely many nonconstant \(f_1,\ldots,f_k\): each nonconstant
polynomial has an irreducible factor, by induction on the degree. Adjoin a
root of an irreducible factor of \(f_1\) to \(E\), then of an irreducible
factor of \(f_2\) over the resulting field, and so on; after \(k\) steps
all \(f_j\) have roots.

Item 2. The evaluation \(E[X]\to K\), \(X\mapsto\alpha\), has nonzero kernel
because \(\alpha\) is algebraic; the kernel is prime, since its image lies in
a field, so it is generated by an irreducible polynomial, the minimal
polynomial \(g\). The image \(E[\alpha]\) is isomorphic to \(E[X]/(g)\), a
field by item~1, so \(E[\alpha]=E(\alpha)\). The map
\(\varepsilon:E[X]\to L\), \(\sum e_iX^i\mapsto\sum\sigma(e_i)\beta^i\), is a
ring homomorphism with \(g\) in its kernel, because \(\sigma g(\beta)=0\).
Its kernel is a proper ideal containing the maximal ideal \((g)\), hence
equal to it, and \(\varepsilon\) induces an embedding
\(E[X]/(g)\to L\). Composing with the isomorphism
\(E(\alpha)\cong E[X]/(g)\) gives
\(\tau(\sum e_i\alpha^i)=\sum\sigma(e_i)\beta^i\), which extends \(\sigma\)
and sends \(\alpha\) to \(\beta\). If \(L\) is algebraically closed, then
\(\sigma g\) has the same positive degree as \(g\), since \(\sigma\) is
injective on the leading coefficient; dividing by its leading coefficient
gives a monic polynomial, which has a root in \(L\).

Item 3. Since \(\alpha\) is transcendental, evaluation at \(\alpha\) is
injective on \(E[X]\), so each element of \(E[\alpha]\) is
\(\sum e_i\alpha^i\) for unique coefficients. The map
\(\sum e_i\alpha^i\mapsto\sum\sigma(e_i)\beta^i\) is therefore a well-defined
ring homomorphism \(E[\alpha]\to L\); it is injective because
\(\sum\sigma(e_i)\beta^i=0\) forces every \(\sigma(e_i)=0\), \(\beta\) being
transcendental over \(\sigma(E)\), and so every \(e_i=0\). The field
\(E(\alpha)\) is the set of quotients of elements of the integral domain
\(E[\alpha]\), and item~4 extends the map to it.

Item 4. Here the fraction field of a subring \(A\) of a field is the set of
quotients \(ab^{-1}\) with \(a,b\in A\), \(b\ne0\), the subfield
generated by \(A\). Let \(\rho:A\to L\) be an injective ring homomorphism.
Put \(\hat\rho(ab^{-1})=\rho(a)\rho(b)^{-1}\), defined because
\(\rho(b)\ne0\). If \(ab^{-1}=a'b'^{-1}\), then \(ab'=a'b\), so
\(\rho(a)\rho(b')=\rho(a')\rho(b)\) and the two values agree. The usual
formulas for sums and products of quotients show that \(\hat\rho\) is a ring
homomorphism, injective because its domain is a field and
\(\hat\rho(1)=1\). Any extension \(\mu\) satisfies
\(\mu(ab^{-1})\rho(b)=\rho(a)\), so it equals \(\hat\rho\).

Item 5. If \(K=\{k_1,\ldots,k_n\}\) were finite, the monic polynomial
\(\prod_i(X-k_i)+1\) of degree \(n\ge1\) would take the value \(1\ne0\) at
every element of \(K\) and so would have no root.
\end{proof}
\end{annexitem}

\begin{annexitem}{annex:p03-w03-p03-020-closure-exists}{thm:p03-w03-p03-020-closure-exists}{Algebraically closed extensions exist}
\begin{proof}
\emph{One step.} Let \(F\) be a field, \(P\) its set of nonconstant
polynomials, and \(R=F[X_f:f\in P]\) the polynomial ring in one variable
for each \(f\in P\). Let \(I\) be the ideal generated by the elements
\(f(X_f)\). If \(I=R\), then \(1=\sum_{j\le k}g_jf_j(X_{f_j})\) for finitely
many \(g_j\in R\). By \cref{thm:p03-w03-p03-020-field-algebra}, item~1,
there is a field \(F'\supseteq F\) in which each \(f_j\) has a root
\(r_j\). The ring homomorphism \(R\to F'\) that is the identity on \(F\),
sends \(X_{f_j}\) to \(r_j\) and every other variable to \(0\), maps the
right side to \(0\) and the left side to \(1\), a contradiction. So \(I\) is
proper. The proper ideals containing \(I\), ordered by inclusion, form a
nonempty family in which every chain has an upper bound, its union, which
is proper because it does not contain \(1\). By Zorn's lemma there is a
maximal ideal \(\mathfrak m\supseteq I\), and \(F^+=R/\mathfrak m\) is a
field. As in item~1 of the lemma, \(F\) embeds in \(F^+\), and every
\(f\in P\) has the root \(X_f+\mathfrak m\).

\emph{Iteration.} Put \(F_0=K\) and \(F_{n+1}=F_n^+\), renamed so that
\(F_n\subseteq F_{n+1}\) as subfields. The union \(\Omega\) is a field. A
nonconstant polynomial over \(\Omega\) has its finitely many coefficients
in some \(F_n\), so it has a root in \(F_{n+1}\subseteq\Omega\). Hence
\(\Omega\) is algebraically closed. It has the characteristic of its
subfield \(K\), since each \(\chi_q\) is an atomic sentence evaluated in
the common prime subfield.

\emph{Algebraic part.} If \(\alpha\) is algebraic over a field \(E\) with
minimal polynomial of degree \(m\), then \(E(\alpha)=E[\alpha]\) is spanned
over \(E\) by \(1,\alpha,\ldots,\alpha^{m-1}\), by the isomorphism with
\(E[X]/(g)\) in the proof of item~2 of the lemma. If
\(\gamma_1,\ldots,\gamma_r\) are algebraic over \(K\), then each is
algebraic over \(K(\gamma_1,\ldots,\gamma_{i-1})\), so by induction
\(K(\gamma_1,\ldots,\gamma_r)\) is spanned over \(K\) by finitely many
products of powers, say \(N\) of them. For any \(\delta\) in this field, the
\(N+1\) elements \(1,\delta,\ldots,\delta^N\) are linearly dependent over
\(K\), so \(\delta\) is algebraic over \(K\). Applied with \(r=2\), this
shows that the set \(\bar K\) of elements of \(\Omega\) algebraic over
\(K\) is closed under the field operations, so it is a subfield. If a
nonconstant \(f\in\bar K[X]\) has coefficients \(\gamma_1,\ldots,\gamma_r\)
and \(\omega\in\Omega\) is a root, then \(\omega\) is algebraic over
\(K(\gamma_1,\ldots,\gamma_r)\), and the same argument applied to
\(\gamma_1,\ldots,\gamma_r,\omega\) shows that \(\omega\) is algebraic over
\(K\). So \(\omega\in\bar K\), and \(\bar K\) is algebraically closed.
\end{proof}
\end{annexitem}

\begin{annexitem}{annex:p03-w03-p03-020-acf-qe}{thm:p03-w03-p03-020-acf-qe}{Elimination for algebraically closed fields of fixed characteristic}
\begin{proof}
We verify the test of \cref{thm:p03-w03-p03-020-one-point}. Let
\(\mathcal M,\mathcal N\entails\mathsf{ACF}_p\), let \(\bar a\in M^n\) and
\(\bar b\in N^n\) have the same quantifier-free type, and let \(c\in M\).

\emph{The generated subrings.} Let \(A=\{t^{\mathcal M}(\bar a):t\text{ an }
\Lang_{\mathrm r}\text{-term in }\bar x\}\), the subring of \(M\)
generated by \(\bar a\), and define
\(\rho(t^{\mathcal M}(\bar a))=t^{\mathcal N}(\bar b)\). Since
\(t^{\mathcal M}(\bar a)=s^{\mathcal M}(\bar a)\) iff the atom \(t=s\)
belongs to the common type iff \(t^{\mathcal N}(\bar b)=s^{\mathcal N}(\bar b)\),
\(\rho\) is well defined and injective. It is a ring homomorphism because
\(\rho\) of a sum or product of term values is the value of the sum or
product term, and it sends each \(a_i\) to \(b_i\). Let \(E\subseteq M\) be
the subfield generated by \(A\). By item~4 of
\cref{thm:p03-w03-p03-020-field-algebra}, \(\rho\) extends to an embedding
\(\sigma:E\to N\).

\emph{Algebraic case.} If \(c\) is algebraic over \(E\), item~2 of the lemma,
with \(N\) algebraically closed, gives an embedding \(\tau:E(c)\to N\)
extending \(\sigma\). Put \(\mathcal N'=\mathcal N\) and \(d=\tau(c)\).

\emph{Transcendental case.} Otherwise, let \(e\) be a new constant and
\[
 \Sigma=\operatorname{ElDiag}(\mathcal N)\cup
 \{q(e)\ne0: q\text{ a nonzero polynomial over }\sigma(E)\},
\]
where the coefficients of \(q\) are written with their parameter constants.
For finitely many nonzero \(q_1,\ldots,q_k\), the product
\(\prod_jq_j\) is a nonzero polynomial over the field \(N\) and has at most
its degree many roots in \(N\); \(N\) is infinite by item~5 of the lemma, so
some element of \(N\) is a root of none of them, and interpreting \(e\) by
it satisfies that finite part of \(\Sigma\). By compactness,
\cref{thm:diagram-tests} and renaming, there are
\(\mathcal N\preccurlyeq\mathcal N'\) and \(d\in N'\) that is a root of no
nonzero polynomial over \(\sigma(E)\), that is, \(d\) is transcendental over
\(\sigma(E)\). Item~3 of the lemma gives an embedding
\(\tau:E(c)\to N'\) extending \(\sigma\) with \(\tau(c)=d\).

\emph{Matching types.} In both cases \(\mathcal N\preccurlyeq\mathcal N'\),
and \(\tau:E(c)\to N'\) is an injective ring homomorphism with
\(\tau(a_i)=b_i\) and \(\tau(c)=d\). By induction on terms, for every
\(\Lang_{\mathrm r}\)-term \(t(\bar x,y)\) the value
\(t^{\mathcal M}(\bar a,c)\) lies in \(E(c)\) and
\(\tau(t^{\mathcal M}(\bar a,c))=t^{\mathcal N'}(\bar b,d)\). Since
\(\tau\) is injective, an atom \(t=s\) holds at \((\bar a,c)\) in
\(\mathcal M\) iff it holds at \((\bar b,d)\) in \(\mathcal N'\). So the two
tuples have the same quantifier-free type, and \(\mathsf{ACF}_p\) eliminates
quantifiers.

\emph{Completeness and model completeness.} The atomic
\(\Lang_{\mathrm r}\)-sentences are equations between closed terms, each
equivalent modulo the field axioms to \(k\cdot1=l\cdot1\) for integers
\(k,l\). In characteristic \(p\) prime this holds iff \(p\) divides \(k-l\),
and in characteristic \(0\) iff \(k=l\). So any two models of
\(\mathsf{ACF}_p\) satisfy the same atomic sentences. The theory is
consistent by \cref{thm:p03-w03-p03-020-closure-exists}, applied to the
field \(\mathbb Z/p\mathbb Z\) or \(\mathbb Q\). Items~5 and~3 of
\cref{thm:p03-w03-p03-020-qe-consequences} give completeness and model
completeness. The characteristic was not used in the elimination; it is
used only to settle the atomic sentences.
\end{proof}
\end{annexitem}

\begin{annexitem}{annex:p03-w03-p03-020-acf-completion}{thm:p03-w03-p03-020-acf-completion}{\(\mathsf{ACF}_p\) is the model completion of \(\mathsf{Fld}_p\)}
\begin{proof}
Item 1. A field \(K\) of characteristic \(p\) is a subfield of an
algebraically closed field \(\Omega\) of the same characteristic by
\cref{thm:p03-w03-p03-020-closure-exists}; \(\Omega\) is a model of
\(\mathsf{ACF}_p\), and the inclusion is an embedding. Every model of
\(\mathsf{ACF}_p\) is a model of \(\mathsf{Fld}_p\), which embeds in itself.
\Cref{thm:p03-w03-p03-020-companion-embedding} gives equality of the
universal consequences.

Item 2 is part of \cref{thm:p03-w03-p03-020-acf-qe}.

Item 3. Consistency: the expansion \(\Omega_K\) with each \(c_k\)
interpreted as \(k\) satisfies \(\operatorname{Diag}(K)\), because the
inclusion is an embedding. Completeness: item~4 of
\cref{thm:p03-w03-p03-020-qe-consequences} with
\(\mathcal A=K\). By \cref{thm:p03-w03-p03-020-qe-completion},
\(\mathsf{ACF}_p\) is the model completion of \(\mathsf{Fld}_p\).

Item 4. The axioms of \(\mathsf{Fld}_p\) are universal sentences, the
quantifier-free sentences \(\chi_p\) or \(\neg\chi_q\), and the sentence
\(\forall x\,\exists y\,(x=0\lor x\cdot y=1)\), so each is of the form
\(\forall\bar x\,\exists\bar y\,\theta\) and \(\mathsf{Fld}_p\) has the same
models as its \(\forall\exists\)-consequences. Item~2 of
\cref{thm:p03-w03-p03-020-companion-ec} identifies the models of
\(\mathsf{ACF}_p\) with the existentially closed models of \(\mathsf{Fld}_p\).
Directly: if \(K\) is not algebraically closed, some monic \(f\) over \(K\)
has no root in \(K\) but has one in an algebraically closed
\(\Omega\supseteq K\), so the existential \(\Lang_{\mathrm r}(K)\)-sentence
\(\exists y\,f(y)=0\) is true in \(\Omega\) and false in \(K\).
\end{proof}
\end{annexitem}

\begin{annexitem}{annex:p03-w03-p03-020-order-extension}{thm:p03-w03-p03-020-order-extension}{Extending an order}
\begin{proof}
\emph{Signs in an ordered field.} In an ordered field, \(0<y\) gives
\(0<y\mathbin{\cdot}y\), and \(y<0\) gives \(0<-y\) by adding \(-y\), so
again \(0<(-y)(-y)=y\mathbin{\cdot}y\); hence every square is \(\ge0\), and
\(0<1\) because \(1=1\cdot1\ne0\). If \(0\le x\) and \(0\le z\) then
\(z\le x+z\) by adding \(z\), so \(0\le x+z\); and \(0<p\), \(0\le w\) give
\(0\le pw\).

\emph{Necessity.} If \(<_L\) extends \(<\), each \(p_i\) with \(0<p_i\) in
\(K\) is positive in \(L\), each \(p_iy_i^2\) is \(\ge0\), and so is their
sum, while \(-1<0\). So no such representation of \(-1\) exists.

\emph{Sufficiency.} Assume there is no such representation. Call
\(P\subseteq L\) a \emph{cone} if \(P+P\subseteq P\), \(P\cdot P\subseteq
P\), \(y^2\in P\) for every \(y\in L\), \(-1\notin P\), and \(P\) contains
every \(p\in K\) with \(0<p\). Let \(S\) be the set of sums
\(\sum_{i\le m}p_iy_i^2\) with \(m\ge1\), \(0<p_i\in K\) and \(y_i\in L\).
It is closed under sums (concatenate the lists) and under products, since
\(\bigl(\sum_ip_iy_i^2\bigr)\bigl(\sum_jq_jz_j^2\bigr)=\sum_{i,j}(p_iq_j)(y_iz_j)^2\)
and \(0<p_iq_j\). It contains each square \(1\cdot y^2\), each positive
\(p=p\cdot1^2\) of \(K\), and it does not contain \(-1\) by assumption. So
\(S\) is a cone. The cones containing \(S\), ordered by inclusion, form a
nonempty family in which the union of a nonempty chain is again a cone,
because any two elements of the union lie in one member of the chain. By
Zorn's lemma there is a maximal cone \(P\supseteq S\).

\emph{\(P\cup-P=L\).} Let \(a\notin P\) and put
\(P'=\{x-aw:x,w\in P\}\). It contains \(P\) (take \(w=0=0^2\)) and \(-a\)
(take \(x=0\), \(w=1\)). It is closed under sums, and under products since
\((x-aw)(x'-aw')=(xx'+a^2ww')-a(xw'+wx')\) with \(a^2ww'\in P\). It
contains all squares and the positive elements of \(K\), because \(P\) does.
If \(-1=x-aw\) with \(x,w\in P\), then \(w\ne0\), since otherwise
\(-1=x\in P\); so \(aw=1+x\) and
\(a=(1+x)\,w\,(w^{-1})^2\in P\), a contradiction. Thus \(P'\) is a cone
containing \(P\), so \(P'=P\) by maximality, and \(-a\in P\).

\emph{\(P\cap-P=\{0\}\).} If \(a\ne0\) and \(a,-a\in P\), then
\(-1=(-a)\,a\,(a^{-1})^2\in P\), which is impossible.

\emph{The order.} Put \(x<_Ly\) iff \(y-x\in P\setminus\{0\}\). It is
irreflexive; it is transitive because \(P+P\subseteq P\) and a sum
\(u+v\) of elements \(u,v\in P\setminus\{0\}\) is not \(0\) (otherwise
\(u=-v\in P\cap-P\)); and for \(x\ne y\) exactly one of \(y-x\) and
\(x-y\) lies in \(P\), so \(<_L\) is a strict linear order. Adding \(z\)
does not change \(y-x\), and if \(0<_Lx\) and \(0<_Ly\) then
\(xy\in P\) and \(xy\ne0\) in a field. So \((L,<_L)\) is an ordered field.
If \(a<b\) in \(K\), then \(b-a\) is a positive element of \(K\), so it lies
in \(P\) and is nonzero, and \(a<_Lb\). Since \(<\) and the restriction of
\(<_L\) to \(K\) are strict linear orders with \(<\) contained in the
restriction, they are equal: if \(a<_Lb\) and not \(a<b\), then \(a=b\) or
\(b<a\), and \(b<a\) gives \(b<_La\); both contradict \(a<_Lb\). So \(<_L\)
extends \(<\).
\end{proof}
\end{annexitem}

\begin{annexitem}{annex:p03-w03-p03-020-root-ext}{thm:p03-w03-p03-020-root-ext}{Adjoined roots keep an order}
\begin{proof}
Both items use \cref{thm:p03-w03-p03-020-order-extension}, and the fact,
from the proof of item~2 of \cref{thm:p03-w03-p03-020-field-algebra}, that
if \(\alpha\) is a root of an irreducible \(g\in K[X]\) of degree \(d\),
then \(g\) is a nonzero constant multiple of the minimal polynomial of
\(\alpha\), \(K(\alpha)=K[\alpha]\cong K[X]/(g)\), and every element of
\(K(\alpha)\) is \(h(\alpha)\) for a unique \(h\in K[X]\) of degree less than
\(d\) or \(h=0\). A polynomial over \(K\) vanishing at \(\alpha\) is
divisible by \(g\).

Item 1. If \(a=b^2\) with \(b\in K\), then \((\alpha-b)(\alpha+b)=0\), so
\(\alpha=\pm b\in K\) and \(K(\alpha)=K\). Otherwise \(X^2-a\) has no root
in \(K\), so it is irreducible of degree \(2\), and each element of
\(K(\alpha)\) is \(x+y\alpha\) for unique \(x,y\in K\). Suppose
\(-1=\sum_ip_i(x_i+y_i\alpha)^2\) with \(0<p_i\in K\). Expanding with
\(\alpha^2=a\),
\[
-1=\sum_ip_i\bigl(x_i^2+ay_i^2\bigr)+\Bigl(2\sum_ip_ix_iy_i\Bigr)\alpha .
\]
By uniqueness, \(-1=\sum_ip_ix_i^2+\sum_i(p_ia)y_i^2\) in \(K\). Every term
on the right is \(\ge0\) in \((K,<)\), because \(0<p_i\) and \(0<p_ia\), so
the sum is \(\ge0>-1\); a contradiction. By
\cref{thm:p03-w03-p03-020-order-extension}, \(<\) extends to
\(K(\alpha)\).

Item 2. Fix \((K,<)\). We prove by strong induction on odd \(d\ge1\) that
the conclusion holds for every irreducible \(g\in K[X]\) of degree \(d\) and
every root \(\alpha\) of \(g\) in any field extending \(K\). For \(d=1\),
\(\alpha\in K\) and \(K(\alpha)=K\). Let \(d\ge3\) be odd, and suppose the
order does not extend to \(K(\alpha)\). By
\cref{thm:p03-w03-p03-020-order-extension} there are \(0<p_i\in K\) and
\(y_i\in K(\alpha)\), \(i\le m\), with \(-1=\sum_ip_iy_i^2\). Write
\(y_i=g_i(\alpha)\) with \(g_i\in K[X]\) zero or of degree at most
\(d-1\), and put
\[
G=1+\sum_ip_ig_i^2\in K[X].
\]
Then \(G(\alpha)=1+(-1)=0\). Let \(e\) be the largest degree of a nonzero
\(g_i\). Not all \(g_i\) are zero, since otherwise \(-1=0\). The coefficient
of \(X^{2e}\) in \(\sum_ip_ig_i^2\) is the sum of
\(p_i\operatorname{lc}(g_i)^2\) over the \(i\) with \(\deg g_i=e\), a sum
of positive elements, hence positive. If \(e=0\), then \(G\) is the
constant \(1+\sum_ip_ig_i^2\ge1\), contradicting \(G(\alpha)=0\). So
\(1\le e\le d-1\), and \(G\) is nonzero of even degree \(2e\le2d-2\).
Since \(G(\alpha)=0\), \(g\) divides \(G\): \(G=gh\) with \(h\ne0\) and
\(\deg h=2e-d\), which is odd and at most \(d-2\). Factor \(h\) into
irreducible polynomials over \(K\) (possible by induction on the degree);
their degrees add up to the odd number \(\deg h\), so some irreducible
factor \(h_1\) has odd degree \(d_1\le d-2\). By item~1 of
\cref{thm:p03-w03-p03-020-field-algebra}, \(K_1=K[X]/(h_1)\) is a field
extending \(K\) in which the class \(\beta\) of \(X\) is a root of \(h_1\),
and \(K_1=K(\beta)\). Then \(G(\beta)=g(\beta)h(\beta)=0\), that is,
\(-1=\sum_ip_ig_i(\beta)^2\) in \(K(\beta)\). By the induction hypothesis
for \(d_1<d\), the order of \(K\) extends to \(K(\beta)\), and the necessity
half of \cref{thm:p03-w03-p03-020-order-extension} forbids this
representation of \(-1\). This contradiction shows that the order extends to
\(K(\alpha)\).
\end{proof}
\end{annexitem}

\begin{annexitem}{annex:p03-w03-p03-020-rcf-char}{thm:p03-w03-p03-020-rcf-char}{Real closed fields by roots}
\begin{proof}
Throughout, \((R,<)\) is an ordered field. As shown at the start of the
proof of \cref{thm:p03-w03-p03-020-order-extension}, squares are \(\ge0\)
and \(0<1\). Hence \(0<n\cdot1\) for every integer \(n\ge1\), so \(R\) has
characteristic \(0\), \(2\) is invertible with \(0<1/2\), and \(R\) is
infinite. For \(x\in R\), \(\abs x\) is \(x\) or \(-x\), whichever is
\(\ge0\). For \(0\le u,v\): if \(u<v\) then \(u^2\le uv<v^2\), so
\(u^2=v^2\) implies \(u=v\), and \(u^2\le v^2\) implies \(u\le v\).
Polynomial functions are those of \cref{def:p03-w03-p03-037-rcf}: the
intermediate-value axiom for \(n\) and coefficients \(a_0,\ldots,a_n\)
concerns the polynomial \(\sum_ka_kX^k\in R[X]\), which may have degree
less than \(n\).

\emph{1 implies 2.} Let \(R\models\mathsf{RCF}\) and \(0<a\). The
polynomial \(X^2-a\) takes the value \(-a<0\) at \(0\) and
\((1+a)^2-a=1+a+a^2>0\) at \(1+a\), and \(0<1+a\); the intermediate-value
axiom with \(n=2\) gives a root. Let \(f\in R[X]\) have odd degree \(n\);
dividing by its leading coefficient we may take \(f\) monic. By item~3 of
\cref{thm:p03-w03-p03-037-poly-basics}, for \(x=M_f+1\) the sign of \(f(x)\)
is \(1\), and for \(x=-(M_f+1)\) it is \((-1)^n=-1\). Since
\(-(M_f+1)<M_f+1\), the intermediate-value axiom gives a root of \(f\).

\emph{2 implies 1.} Assume 2.

\emph{Step 1: the field \(R(\mathrm i)\).} Since \(-1<0\) and squares are
\(\ge0\), \(-1\) is not a square, so \(X^2+1\) has no root in \(R\) and is
irreducible. By item~1 of \cref{thm:p03-w03-p03-020-field-algebra},
\(R(\mathrm i)=R[X]/(X^2+1)\) is a field whose elements are
\(a+b\mathrm i\) with unique \(a,b\in R\), and \(\mathrm i^2=-1\). Item~2 of
the same lemma, applied to the identity of \(R\) and the root
\(-\mathrm i\) of \(X^2+1\), gives a ring homomorphism
\(\kappa:R(\mathrm i)\to R(\mathrm i)\) with
\(\kappa(a+b\mathrm i)=a-b\mathrm i\); it is its own inverse, so it is an
automorphism, and its fixed elements are exactly those of \(R\). Applied to
coefficients, \(\kappa\) is a ring automorphism of \(R(\mathrm i)[Y]\)
whose fixed polynomials are those of \(R[Y]\); write \(\bar f\) for
\(\kappa f\).

\emph{Step 2: every element of \(R(\mathrm i)\) is a square.} By 2, every
\(x\ge0\) in \(R\) has a square root; one of \(\pm\) it is \(\ge0\), and we
write \(\sqrt x\) for that one, which is unique by the remark on
\(u^2=v^2\). Let \(a+b\mathrm i\in R(\mathrm i)\) and put
\(r=\sqrt{a^2+b^2}\). Then \(r^2\ge a^2=\abs a^2\), so \(r\ge\abs a\) and
\(r+a\ge0\), \(r-a\ge0\). Put \(c=\sqrt{(r+a)/2}\) and
\(d_0=\sqrt{(r-a)/2}\), and let \(d=d_0\) if \(b\ge0\) and \(d=-d_0\) if
\(b<0\). Then \(c^2-d^2=a\) and \((2cd)^2=4c^2d_0^2=(r+a)(r-a)=r^2-a^2=b^2\).
Moreover \(2cd\) is \(\ge0\) when \(b\ge0\) and \(\le0\) when \(b<0\), so
\(\abs{2cd}\) and \(\abs b\) are nonnegative with equal squares, hence equal,
and \(2cd=b\). Therefore
\((c+d\mathrm i)^2=c^2-d^2+2cd\,\mathrm i=a+b\mathrm i\).

\emph{Step 3: quadratics over \(R(\mathrm i)\).} Let
\(s,q\in R(\mathrm i)\) and choose \(w\) with \(w^2=s^2-4q\) by Step~2.
Then \(z=(-s+w)/2\) satisfies
\(z^2+sz+q=\bigl(s^2-2sw+w^2-2s^2+2sw+4q\bigr)/4=(w^2-s^2+4q)/4=0\).

\emph{Step 4: symmetric polynomials.} Let \(A\) be a commutative ring and
\(e_1,\ldots,e_n\in A[X_1,\ldots,X_n]\) the elementary symmetric
polynomials, defined by
\(\prod_{j\le n}(Y-X_j)=\sum_{k\le n}(-1)^ke_kY^{n-k}\) with \(e_0=1\).
\emph{Claim:} if \(s\in A[X_1,\ldots,X_n]\) is unchanged by every
permutation of the variables and has total degree at most \(D\) (or
\(s=0\)), then \(s=P(e_1,\ldots,e_n)\) for some
\(P\in A[Z_1,\ldots,Z_n]\) all of whose monomials
\(Z_1^{m_1}\cdots Z_n^{m_n}\) have weight \(m_1+2m_2+\cdots+nm_n\le D\).
Since \(e_k\) has total degree \(k\), such a \(P(e_1,\ldots,e_n)\) has total
degree at most \(D\). The proof is by induction on \(n\) and, for fixed
\(n\), on \(D\). For \(n=1\), \(e_1=X_1\) and \(P=s\). Let \(n\ge2\). For
\(D=0\), \(s\) is a constant. Let \(D\ge1\). Setting \(X_n=0\) gives a
polynomial \(s_0(X_1,\ldots,X_{n-1})\) unchanged by permutations of
\(X_1,\ldots,X_{n-1}\), of total degree at most \(D\); and
\(e_k(X_1,\ldots,X_{n-1},0)\) is the \(k\)-th elementary symmetric
polynomial \(e'_k\) of \(X_1,\ldots,X_{n-1}\) for \(k<n\), while
\(e_n(X_1,\ldots,X_{n-1},0)=0\). By the induction hypothesis on \(n\),
\(s_0=Q(e'_1,\ldots,e'_{n-1})\) with \(Q\) of weight at most \(D\). Put
\(u=s-Q(e_1,\ldots,e_{n-1})\). It is symmetric, of total degree at most
\(D\), and \(u(X_1,\ldots,X_{n-1},0)=0\), so every monomial of \(u\)
contains \(X_n\). If a monomial of \(u\) did not contain \(X_j\), the
transposition of \(X_j\) and \(X_n\), which leaves \(u\) unchanged, would
send it to a monomial of \(u\), with the same nonzero coefficient, not
containing \(X_n\). So every monomial of \(u\) is divisible by
\(X_1\cdots X_n=e_n\), and \(u=e_nv\) with \(v=0\) or of total degree at
most \(D-n\). Multiplication by the monomial \(e_n\) is injective on
\(A[X_1,\ldots,X_n]\), since it shifts exponents, and \(e_n\) is symmetric;
so for each permutation \(\pi\), \(e_n\,\pi v=\pi u=u=e_nv\) gives
\(\pi v=v\). By the induction hypothesis on \(D\),
\(v=P'(e_1,\ldots,e_n)\) with \(P'\) of weight at most \(D-n\), and
\(s=Q(e_1,\ldots,e_{n-1})+e_nP'(e_1,\ldots,e_n)\) has the required form.

\emph{Step 5: every nonconstant \(f\in R[Y]\) has a root in
\(R(\mathrm i)\).} Write \(\deg f=n=2^km\) with \(m\) odd; we induct on
\(k\). If \(k=0\), \(f\) has a root in \(R\) by 2. Let \(k\ge1\). Dividing by
the leading coefficient, take \(f\) monic. Adjoining roots one at a time by
item~1 of \cref{thm:p03-w03-p03-020-field-algebra} and dividing out the
corresponding linear factors, we obtain a field \(E\supseteq R(\mathrm i)\)
with \(f=\prod_{j\le n}(Y-\alpha_j)\) in \(E[Y]\). Comparing with Step~4,
\(e_k(\alpha_1,\ldots,\alpha_n)=(-1)^k\cdot\)(the coefficient of
\(Y^{n-k}\) in \(f\)) lies in \(R\). For \(c\in R\) let
\[
H_c=\prod_{1\le j<l\le n}\bigl(Y-(X_j+X_l+cX_jX_l)\bigr)\in
R[X_1,\ldots,X_n][Y].
\]
A permutation \(\pi\) of the variables permutes the factors, since the
factor for \(\{j,l\}\) is unchanged when \(j\) and \(l\) are exchanged; so
each coefficient of a power of \(Y\) in \(H_c\) is a symmetric polynomial
over \(R\), hence by Step~4 a polynomial over \(R\) in \(e_1,\ldots,e_n\).
Evaluating the \(X_j\) at the \(\alpha_j\) therefore gives a monic
\(G_c=\prod_{j<l}\bigl(Y-\beta_{jl}(c)\bigr)\in R[Y]\), where
\(\beta_{jl}(c)=\alpha_j+\alpha_l+c\alpha_j\alpha_l\). Its degree is
\(N=n(n-1)/2=2^{k-1}m(2^km-1)\), and \(m(2^km-1)\) is odd because
\(k\ge1\). By the induction hypothesis \(G_c\) has a root \(z\in
R(\mathrm i)\); then \(\prod_{j<l}(z-\beta_{jl}(c))=0\) in the field \(E\),
so \(z=\beta_{jl}(c)\) for some pair, and \(\beta_{jl}(c)\in R(\mathrm i)\).
Since \(R\) is infinite, take \(N+1\) distinct values of \(c\); two of
them, \(c\ne c'\), give the same pair \(\{j,l\}\). With
\(\sigma=\alpha_j+\alpha_l\) and \(\rho=\alpha_j\alpha_l\), both
\(\sigma+c\rho\) and \(\sigma+c'\rho\) lie in \(R(\mathrm i)\), so
\(\rho=\bigl((\sigma+c\rho)-(\sigma+c'\rho)\bigr)/(c-c')\in R(\mathrm i)\)
and then \(\sigma\in R(\mathrm i)\). By Step~3 the polynomial
\(Y^2-\sigma Y+\rho=(Y-\alpha_j)(Y-\alpha_l)\) has a root
\(z'\in R(\mathrm i)\), and \((z'-\alpha_j)(z'-\alpha_l)=0\) in \(E\) gives
\(z'\in\{\alpha_j,\alpha_l\}\). So \(f\) has the root \(z'\) in
\(R(\mathrm i)\).

\emph{Step 6: \(R(\mathrm i)\) is algebraically closed.} Let
\(f\in R(\mathrm i)[Y]\) be nonconstant and put \(h=f\bar f\). Then
\(\bar h=\bar ff=h\), so \(h\in R[Y]\), and \(h\) is nonconstant because
\(\deg\bar f=\deg f\). By Step~5, \(h(z)=0\) for some
\(z\in R(\mathrm i)\), so \(f(z)=0\) or \(\bar f(z)=0\); in the second case
applying \(\kappa\) gives \(f(\kappa z)=0\). Either way \(f\) has a root in
\(R(\mathrm i)\).

\emph{Step 7: irreducible polynomials over \(R\).} Let \(g\in R[X]\) be
monic and irreducible. By Step~6 it has a root \(z=a+b\mathrm i\in
R(\mathrm i)\). Then \((z-a)^2=-b^2\), so \(z\) is a root of
\(Q=(X-a)^2+b^2\in R[X]\). Since \(g\) is a constant multiple of the
minimal polynomial of \(z\) over \(R\), \(g\) divides \(Q\), and
\(\deg g\le2\). If \(\deg g=2\), write \(g=X^2+dX+e\). If
\(d^2/4-e\ge0\), then with \(w=\sqrt{d^2/4-e}\) the element \(-d/2+w\) of
\(R\) would be a root of \(g\), as the computation of Step~3 shows; so
\(e-d^2/4>0\), and \(g(x)=(x+d/2)^2+(e-d^2/4)>0\) for every \(x\in R\).

\emph{Step 8: the intermediate-value axioms.} Let \(p\in R[X]\),
\(u<v\) and \(p(u)p(v)<0\); in particular \(p\ne0\). Factor
\(p=a\prod_j(X-r_j)\prod_kg_k\) with \(a\in R\setminus\{0\}\), \(r_j\in R\)
and \(g_k\) monic irreducible of degree \(2\), using Step~7 and induction
on the degree. Then
\[
p(u)p(v)=a^2\prod_j(u-r_j)(v-r_j)\prod_kg_k(u)g_k(v),
\]
where \(a^2>0\) and each \(g_k(u)g_k(v)>0\) by Step~7. So some
\((u-r_j)(v-r_j)\) is negative, and since \(u-r_j<v-r_j\) this means
\(u-r_j<0<v-r_j\). Thus \(u<r_j<v\) and \(p(r_j)=0\). This is the
intermediate-value axiom of \cref{def:p03-w03-p03-037-rcf} for every \(n\)
and all coefficients, so \((R,<)\models\mathsf{RCF}\).

The final statement of the theorem was shown in Steps~1 and~6.
\end{proof}
\end{annexitem}

\begin{annexitem}{annex:p03-w03-p03-020-real-closure}{thm:p03-w03-p03-020-real-closure}{Real closures exist}
\begin{proof}
By \cref{thm:p03-w03-p03-020-closure-exists}, \(F\) is a subfield of an
algebraically closed field, and the elements algebraic over \(F\) form an
algebraically closed field \(\Omega\) algebraic over \(F\); so a field
\(\Omega\) as in the statement exists. Fix one.

Let \(\mathcal E\) be the set of pairs \((E,<_E)\) with \(E\) a subfield of
\(\Omega\) containing \(F\) and \(<_E\) an order making \((E,<_E)\) an
ordered field that extends \(<\). Order \(\mathcal E\) by
\((E,<_E)\sqsubseteq(E',<_{E'})\) iff \(E\subseteq E'\) and \(<_{E'}\)
extends \(<_E\). The family contains \((F,<)\). If
\(\{(E_\lambda,<_\lambda)\}\) is a nonempty chain, its union \(E^*\) is a
subfield of \(\Omega\), and \(x<^*y\) iff \(x<_\lambda y\) for some
\(\lambda\) with \(x,y\in E_\lambda\) is well defined, since the orders agree
where both are defined. Any finitely many elements lie in one
\(E_\lambda\), so the order and field axioms hold in \((E^*,<^*)\), which
extends each member of the chain and in particular \(<\). By Zorn's lemma
\(\mathcal E\) has a maximal element \((R,<_R)\).

\emph{Positive elements are squares.} Let \(0<_Ra\) and suppose \(a\) is not
a square in \(R\). Since \(\Omega\) is algebraically closed there is
\(\beta\in\Omega\) with \(\beta^2=a\), and \(\beta\notin R\). By item~1 of
\cref{thm:p03-w03-p03-020-root-ext}, \(<_R\) extends to an order of
the subfield \(R(\beta)\) of \(\Omega\). That order extends \(<\), so
\(R(\beta)\) with it is a member of \(\mathcal E\) strictly above
\((R,<_R)\), contradicting maximality.

\emph{Odd-degree polynomials have roots.} Let \(f\in R[X]\) have odd degree.
Its irreducible factors have degrees adding up to \(\deg f\), so one of
them, \(g\), has odd degree. Let \(\alpha\in\Omega\) be a root of \(g\). By
item~2 of \cref{thm:p03-w03-p03-020-root-ext}, \(<_R\) extends to
an order of \(R(\alpha)\subseteq\Omega\), and maximality gives
\(R(\alpha)=R\). So \(\alpha\in R\) and \(f(\alpha)=0\).

By \cref{thm:p03-w03-p03-020-rcf-char},
\((R,<_R)\models\mathsf{RCF}\). Every element of \(R\subseteq\Omega\) is
algebraic over \(F\), and \(<_R\) extends \(<\). So \((R,<_R)\) is a real
closure of \((F,<)\).
\end{proof}
\end{annexitem}

\begin{annexitem}{annex:p03-w03-p03-020-rcf-mc}{thm:p03-w03-p03-020-rcf-mc}{Model completion of ordered fields}
\begin{proof}
Item 1. If \((F,<)\models\mathsf{OF}\), a real closure \((R,<_R)\) exists
by \cref{thm:p03-w03-p03-020-real-closure}, and the inclusion
\(F\subseteq R\) is an \(\Lang_{\mathrm{or}}\)-embedding because \(F\) is a
subfield and \(<\) is the restriction of \(<_R\). Every model of
\(\mathsf{RCF}\) is a model of \(\mathsf{OF}\) and embeds in itself. By
\cref{thm:p03-w03-p03-020-companion-embedding},
\(\mathsf{OF}_\forall=\mathsf{RCF}_\forall\).

Item 2. Elimination is \cref{thm:p03-w03-p03-037-qe}. Model completeness
follows by item~3 of \cref{thm:p03-w03-p03-020-qe-consequences}.
Completeness is \cref{thm:p03-w03-p03-037-complete}; it also follows from
item~5 of \cref{thm:p03-w03-p03-020-qe-consequences}, since every atomic
\(\Lang_{\mathrm{or}}\)-sentence is \(k\cdot1=l\cdot1\) or
\(k\cdot1<l\cdot1\) for integers \(k,l\) after evaluating the closed terms,
and in every ordered field these hold iff \(k=l\), respectively \(k<l\),
because \(0<n\cdot1\) for \(n\ge1\); and \(\mathsf{RCF}\) is consistent by
item~1 applied to \(\mathbb Q\).

Item 3. By items~1 and~2, \(\mathsf{RCF}_\forall=\mathsf{OF}_\forall\) and
\(\mathsf{RCF}\) eliminates quantifiers, so
\cref{thm:p03-w03-p03-020-qe-completion} shows that \(\mathsf{RCF}\) is the
model completion of \(\mathsf{OF}\) and that
\(\mathsf{RCF}\cup\operatorname{Diag}(F,<)\) is consistent and complete for
every ordered field \((F,<)\). If \((R_1,<_1)\) and \((R_2,<_2)\) are real
closures of \((F,<)\), the inclusions of \((F,<)\) are embeddings, so the
expansions interpreting each \(c_a\) as \(a\) satisfy
\(\operatorname{Diag}(F,<)\) by \cref{thm:diagram-tests}. Both are models of
the complete theory \(\mathsf{RCF}\cup\operatorname{Diag}(F,<)\), so they
satisfy the same \(\Lang_{\mathrm{or}}(F)\)-sentences.

Item 4. The axioms of \(\mathsf{OF}\) listed in
\cref{def:p03-w03-p03-020-ordered-fields} are universal sentences together
with \(\forall x\,\exists y\,(x=0\lor x\mathbin{\cdot}y=1)\); each has the
form \(\forall\bar x\,\exists\bar y\,\theta\) with \(\theta\)
quantifier-free, so \(\mathsf{OF}\) and \(\mathsf{OF}_{\forall\exists}\)
have the same models. \Cref{thm:p03-w03-p03-020-companion-ec}, item~2,
applied to the model companion \(\mathsf{RCF}\) of \(\mathsf{OF}\),
identifies the models of \(\mathsf{RCF}\) with the ordered fields that are
existentially closed among ordered fields.

The direct witness for one direction: if \((F,<)\) is an ordered field that
is not real closed, then by \cref{thm:p03-w03-p03-020-rcf-char}
either some \(a\) with \(0<a\) is not a square in \(F\), or some \(f\in
F[X]\) of odd degree has no root in \(F\). In a real closure \((R,<_R)\) of
\((F,<)\), which is an ordered field extending \((F,<)\), \(a\) is a square
and \(f\) has a root by the same theorem. So the existential
\(\Lang_{\mathrm{or}}(F)\)-sentence \(\exists y\,(y\mathbin{\cdot}y=c_a)\),
respectively \(\exists y\,(t_f(y)=c_0)\), where \(t_f\) is a term with
parameters in \(F\) for \(f\), is true in \(R\) and false in \(F\), and
\((F,<)\) is not existentially closed among ordered fields.
\end{proof}
\end{annexitem}
